\documentclass[11pt,a4paper]{article}
\usepackage[utf8]{inputenc}
\usepackage[french]{babel}
\usepackage[T1]{fontenc}
\usepackage{lmodern}
\usepackage[margin=1.5cm]{geometry}

\begin{document}

% En-tête
\begin{center}
    {\Huge \textbf{Léo Martin}} \\[0.2cm]
    {\Large \textbf{Chercheur en Bio-informatique}} \\[0.1cm]
    léo.martin@email.com \ | \ 06 66 51 32 35
\end{center}

\vspace{0.3cm}

% Résumé / Profil
\section*{Profil Professionnel}
\noindent
Bio-informaticien doté d'une double compétence en biologie et en développement informatique. J'accompagne les équipes de recherche dans le traitement, l'analyse et la visualisation de données génomiques à grande échelle pour accélérer les découvertes médicales. Je maîtrise la mise en place de pipelines reproductibles et le déploiement sur des clusters de calcul haute performance (HPC).

\vspace{0.3cm}

% Compétences
\section*{Compétences Techniques}
\noindent
\textbf{Python} $\bullet$ \textbf{Nextflow} $\bullet$ \textbf{NGS} $\bullet$ \textbf{Docker} $\bullet$ \textbf{Bases de données biologiques} $\bullet$ \textbf{Analyse Génomique} $\bullet$ \textbf{BioPython} $\bullet$ \textbf{Statistiques}

\vspace{0.3cm}

% Expériences
\section*{Expériences Professionnelles}
\noindent \textbf{Chercheur en Bio-informatique / Confirmé} --- \textit{Inserm} \hfill \textbf{12/2022 à Aujourd'hui} \\
Création d'applications web interactives avec R Shiny permettant aux biologistes d'explorer visuellement les résultats d'expression génique (RNA-Seq). Gestion et optimisation des requêtes sur un cluster de calcul haute performance (HPC) sous Slurm pour le traitement de pétaoctets de données brutes. Automatisation du nettoyage et du contrôle qualité (QC) des données issues des séquenceurs Illumina, réduisant les erreurs d'alignement de 15\%. Modélisation structurelle 3D de protéines et simulation de dynamique moléculaire, permettant de valider l'affinité de nouvelles molécules thérapeutiques. \\[0.3cm]
\noindent \textbf{Chercheur en Bio-informatique} --- \textit{Institut Pasteur} \hfill \textbf{09/2020 à 03/2022} \\
Automatisation des rapports de diagnostic génétique à l'aide de scripts RMarkdown, garantissant une traçabilité totale pour les audits de santé. Rédaction de rapports d'analyse statistique détaillés avec RMarkdown pour orienter les décisions lors des essais cliniques de phase 2 du laboratoire. Analyse de larges jeux de données de séquençage (NGS) à l'aide de scripts Python et R, aboutissant à l'identification de nouveaux biomarqueurs tumoraux. Croisement et intégration de multiples bases de données biologiques publiques (NCBI, Ensembl, ClinVar) pour annoter des variants génétiques rares. \\[0.3cm]
\noindent \textbf{Stage — Assistant Chercheur en Bio-informatique} --- \textit{Sanofi} \hfill \textbf{05/2018 à 06/2020} \\
Développement de pipelines bio-informatiques reproductibles avec Nextflow et Docker, réduisant le temps d'analyse des génomes de 3 jours à 12 heures. Conception d'une base de données relationnelle sécurisée (HDS) centralisant les données phénotypiques et génotypiques d'une cohorte de 10 000 patients. Développement d'un algorithme de machine learning capable de prédire l'interaction protéine-ligand avec une précision de 88\%. \\[0.3cm]


\vspace{0.3cm}

% Formations
\section*{Formation \& Diplômes}
\noindent \textbf{Licence Professionnelle en Data Science pour la Santé} \hfill \textbf{2019 - 2020} \\
\textit{Université d'Évry} \\[0.3cm]
\noindent \textbf{BTS en Data Science pour la Santé} \hfill \textbf{2017 - 2019} \\
\textit{IUT Informatique option Santé} \\[0.3cm]


\end{document}